% !PRINTED 350
\weberpage{350}
%<<CONT

\[
\begin{aligned}
&\text{a)}\quad x^3 + e x - g = 0\\
&\text{b)}\quad x^3 - e x - g = 0, &&\frac{e^3}{27} < \frac{g^2}{4}\\
&\text{c)}\quad x^3 - e x - g = 0, &&\frac{e^3}{27} > \frac{g^2}{4}.
\end{aligned}
\]

Die Grenzfälle, dass eine der Grössen $e$, $g$, $4e^3 - 27g^2$ verschwindet, schliessen wir aus.

Um die \gesperrt{Cardanische} Formel anzuwenden (\S.\,35), setzen wir

\weq{2}{R = \frac{g^2}{4} \pm \frac{e^3}{27},}

wo das obere Zeichen im Falle a), das untere in den Fällen b) und c) gilt. In den Fällen a) und b) ist nur eine Wurzel reell, in dem Falle c) (dem Casus irreducibilis) sind alle drei Wurzeln reell. Die \gesperrt{Cardanische} Formel giebt

\weq{3}{x = \sqrt[3]{\frac{g}{2} + \sqrt{R}} + \sqrt[3]{\frac{g}{2} - \sqrt{R}}.}

Wir verfahren nun so in den drei Fällen:

a) Wir führen einen Winkel $\vartheta$ ein, den wir zwischen $0^0$ und $90^0$ wählen können, durch die Gleichung:

\weq{4}{\frac{g}{2} = \sqrt{\frac{e^3}{27}}\,\mathrm{cotg}\,\vartheta,}

also

\[
\sqrt{R} = \sqrt{\frac{e^3}{27}}\,\frac{1}{\sin\vartheta},
\]

\[
\begin{aligned}
x &= \sqrt{\frac{e}{3}}\left(\sqrt[3]{\mathrm{cotg}\,\vartheta + \frac{1}{\sin\vartheta}} + \sqrt[3]{\mathrm{cotg}\,\vartheta - \frac{1}{\sin\vartheta}}\right)\\
&= \sqrt{\frac{e}{3}}\left(\sqrt[3]{\mathrm{cotg}\,\frac{\vartheta}{2}} - \sqrt[3]{\mathrm{tg}\,\frac{\vartheta}{2}}\right),
\end{aligned}
\]

und wenn wir noch einen Winkel $\varphi$ aus

\weq{5}{\mathrm{tg}\,\frac{\vartheta}{2} = \mathrm{tg}\,\varphi^3}

bestimmen, so ergiebt sich

\weq{6}{x = 2\sqrt{\frac{e}{3}}\,\mathrm{cotg}\,2\varphi.}

Die Winkel $\vartheta$, $\varphi$ und zuletzt $x$ findet man aus den logarithmisch trigonometrischen Tafeln nach den Formeln (4), (5), (6).
